% ============================================================
% ch23.tex —— 第 23 章 素数定理：素数的交通流量图
% ============================================================
\chapter{素数定理：素数的交通流量图}

第 20 章的预报员 $x/\ln x$ 已经上岗，比值从 $1.15$ 一路收敛到 $1.05$——方向对，但精度不够体面。这一章做两件事：给预报员进修（进修成果 $\mathrm{li}(x)$），然后正式宣读定理全文，并围观人类为了证明它造出的最重的机器。

\section{预报员进修班：从"一口价"到"逐段累积"}

$x/\ln x$ 的毛病在哪？它是\textbf{一口价}：拿 $x$ 这一个点上的密度 $\tfrac1{\ln x}$，乘以全程长度 $x$——相当于用终点附近的物价，估算整段高速路的总花费。而第 22 章末的推导早就指出：密度在\textbf{沿途变化}（$\tfrac1{\ln t}$ 随 $t$ 一路下降），正确算法是\textbf{逐段累积}：
\[
\mathrm{li}(x) \;:=\; \int_2^x \frac{dt}{\ln t} \qquad (\text{正式名：对数积分，logarithmic integral}).
\]
$\mathrm{li}$ 读"利"。它的直觉：从 $2$ 走到 $x$，每前进一小段 $dt$，新增约 $\tfrac{dt}{\ln t}$ 个素数——第 17、18 章的"堆积"在数论里上岗。高斯 1792 年猜的\textbf{正是它}（"密度 $= \tfrac1{\ln x}$"的积分式表述），而不是 $x/\ln x$——后者只是 $\mathrm{li}$ 的一阶速算版。

数值三列并排（体会"一口价"与"逐段累积"的差距）：

\begin{center}
\begin{tabular}{ccccc}
\toprule
$x$ & $\pi(x)$ & $x/\ln x$（误差） & $\mathrm{li}(x)$（误差） \\
\midrule
$10^3$ & $168$ & $144.8$（$-13.8\%$） & $177.6$（$+5.7\%$） \\
$10^4$ & $1\,229$ & $1\,085.7$（$-11.7\%$） & $1\,246.1$（$+1.4\%$） \\
$10^6$ & $78\,498$ & $72\,382$（$-7.8\%$） & $78\,627.5$（$+0.16\%$） \\
$10^9$ & $50\,847\,534$ & $48\,254\,942$（$-5.1\%$） & $50\,849\,235$（$+0.003\%$） \\
$10^{18}$ & $24\,739\,954\,287\,740\,860$ & $2.36\times10^{16}$（$-4.6\%$） & $2.4760\times10^{16}$（$+0.00008\%$） \\
\bottomrule
\end{tabular}
\end{center}

读表三得：其一，$x/\ln x$ \textbf{永远低估}，误差收得很慢（$10^{18}$ 处还差 $4.6\%$）；其二，$\mathrm{li}$ \textbf{永远高估}（在人类算得动的范围内），误差崩塌式收窄（$10^{18}$ 处只剩千万分之一）；其三，两个预报员一低一高，把真值 $\pi(x)$ 夹在中间——\textbf{又见第 14 章的夹逼结构}。

\section{定理宣读}

\begin{dingli}[素数定理（Hadamard 1896；de la Vall\'ee Poussin 1896）]
\[
\lim_{x \to \infty} \frac{\pi(x)}{\mathrm{li}(x)} = 1
\qquad\left(\text{等价地，} \lim_{x\to\infty} \frac{\pi(x)}{x/\ln x} = 1\right).
\]
\end{dingli}

定理的字面意思是"比值落脚于 $1$"（渐近），\textbf{不}承诺误差消失。事实上 $\pi(x) - \mathrm{li}(x)$ 的性态是更深的问题：在人类算过的范围内它恒负（$\mathrm{li}$ 偏高），但李特尔伍德 1914 年证明\textbf{这个差会无限次变号}——只是第一次变号的位置大得骇人：斯丘斯 1955 年给出上界，第一次变号发生在 $10^{10^{10^{964}}}$ 之前（一个写不进宇宙的数，后来被压到 $10^{316}$ 附近量级的讨论仍在继续）。\textbf{素数定理连自己的误差都长得不讲道理}——而治理误差的钥匙，正是第 24 章的零点。

\section{为什么必须动用重型机器}

素数定理证明的困难是结构性的。回顾全书：$\pi(x)$ 是\textbf{台阶函数}——每遇到一个素数跳一格，其余时刻纹丝不动，土得掉渣；$\mathrm{li}(x)$ 是\textbf{光滑曲线}——处处可导，优雅无棱角。\textbf{定理说：最粗粝的台阶，在最粗的尺度上，被最光滑的曲线统治}。让台阶贴上曲线，等价于让"素数的随机抖动"在宏观上互相抵消——而证明"抖动会抵消"，第 22 章的初等机器（欧拉乘积、实数范围的 $\zeta$）功率不够，必须把 $\zeta$ 搬进\textbf{复平面}，用复分析的整体性（解析延拓、零点分布）来摁住抖动。1896 年的证明正是：证出"$\zeta$ 在 $\mathrm{Re}(s) = 1$ 这条线上没有零点" $\Rightarrow$ 素数定理。

这个"离散贴连续"的怪圈值得再停一步体会——它是二十世纪数学反复穿越的秘境：随机矩阵（量子物理）与 $\zeta$ 零点的统计高度吻合（第 24 章）、素数与声学的类比（把素数当"频率"能合成自然的鼓声）……\textbf{离散与连续之间的暗桥不止一座}，素数定理是最早被发现的那座。

\begin{gongju}
\textbf{数值常识补给（心算素数的世界）}$\ln 10 \approx 2.303$；一个 $n$ 位数 $x\approx10^{n-1}$，$\ln x \approx 2.303(n-1)$。于是"$x = 10^k$ 以内素数约 $\dfrac{10^k}{2.303k}$ 个"可心算：$k=9$ 时 $\dfrac{10^9}{20.7} \approx 4.8\times10^7$，与真值 $5.08\times10^7$ 差 $5\%$。"$10^{100}$ 以内素数个数"这种题，现在是你的一行除法。另一个方向：\textbf{第 $n$ 个素数} $p_n \approx n\ln n$（素数定理的逆读）：$n = 10^6$ 时 $p_n \approx 10^6\times13.8 = 1.38\times10^7$，真值 $15\,485\,863$，误差 $11\%$；进修版 $n(\ln n + \ln\ln n - 1) \approx 10^6(13.82+2.63-1) = 1.54\times10^7$，误差 $0.4\%$。
\end{gongju}

\begin{lishi}
素数定理是"猜想一百年才被证明"的标本：欧拉 1737 埋种（密度思想的雏形）；高斯（1792，15 岁）与勒让德（1798）独立提出；切比雪夫（1852）初等方法夹逼到 $0.92$--$1.11$ 倍，离终点一步却止步（初等工具的极限——他证明了：若 $\pi(x)/(x/\ln x)$ 有极限，极限必为 $1$，但"极限存在"卡了四十年）；黎曼 1859 年把 $\zeta$ 搬进复平面（八页论文，本意正是攻克素数定理，没攻克但留下了地图）；1896 年阿达马与德拉瓦莱-普桑\textbf{同年、独立}完成冲线。1948 年塞尔伯格与爱多士又找到不用复分析的"初等证明"——轰动一时，但"初等"只指"不用复分析"，其难度连职业数学家都望而生畏；这提示了一个深刻的事实：\textbf{素数定理的"难度"不在工具的复杂，而在对象本身}。
\end{lishi}

\begin{dongshou}
\begin{itemize}
  \yisi 补齐表格的 $10^5$ 行：$\pi(10^5) = 9592$，算 $x/\ln x = 10^5/11.513 = 8686$（误差 $-9.4\%$）、$\mathrm{li}(10^5) \approx 9629$（误差 $+0.4\%$），确认"$x/\ln x$ 误差缓降、$\mathrm{li}$ 误差崩塌"的走势。
  \yier 用"进修班速成公式"$\mathrm{li}(x) \approx \dfrac{x}{\ln x}\left(1 + \dfrac1{\ln x}\right)$ 重算 $10^6$ 行：$\dfrac{10^6}{13.816}\times1.0724 = 77\,608$（误差从 $-7.8\%$ 砍到 $-1.1\%$）——一阶修正立竿见影。
  \yier 用"第 $n$ 个素数"公式估算：(1) 第 $1000$ 个素数——速算 $n\ln n = 1000\times6.908 = 6908$，真值 $7919$（偏低 $13\%$）；进修版 $n(\ln n + \ln\ln n - 1) = 1000\times(6.908+1.933-1) = 7841$（误差 $1\%$）。(2) 第 $10^7$ 个素数——速算 $10^7\times16.118 = 1.612\times10^8$，真值 $179\,424\,673 = 1.794\times10^8$（偏低 $10\%$）；进修版 $10^7\times(16.118+2.780-1) = 1.790\times10^8$（误差 $0.2\%$）。两问并看："速算—进修—真值"的三明治结构清晰可见。
  \yisi 利用心算补给秒答：$10^{50}$ 以内大约多少素数？（$\dfrac{10^{50}}{2.303\times50} \approx 8.7\times10^{47}$。再换算成"平均几个数一个素数"：约 $115$ 个。）
\end{itemize}
\end{dongshou}

\begin{shaonao}
\textbf{素数定理的三个"等价马甲"。}同一个定理在不同语言里的样子（互相可推）：

(1) 密度版：$x$ 附近的素数密度 $\to \dfrac1{\ln x}$（高斯的原始眼光）；

(2) 计数版：$\pi(x) \sim \dfrac{x}{\ln x}$（勒让德的眼光）；

(3) 间距版：第 $n$ 个素数 $p_n \sim n\ln n$（把计数版反解）。

请完成三个版式的互推练习：(a) 由 (2) 推 (3)：令 $n = \pi(p_n) \approx \dfrac{p_n}{\ln p_n}$，反解 $p_n \approx n\ln p_n$，再"自迭代"一次（$\ln p_n \approx \ln(n\ln n) \approx \ln n$，因为 $\ln\ln n$ 比 $\ln n$ 小一截）得 $p_n \approx n\ln n$。(b) 由 (3) 推 (2)：把 $x = p_n$ 代回。这两个推演让你亲手把"计数"与"间距"缝在一起——\textbf{同一个现象的三件外套，数学家的日常就是互译}。第四件外套"$\sum_{p\le x}\tfrac1p \sim \ln\ln x$"（第 22 章）与它们只差半步等价（Mertens 定理的精确版），留给你日后核对。
\end{shaonao}

\begin{chengshi}
本章全程只陈列、不证明。素数定理的完整证明（哪怕初等版）远超本书工具箱——这是全书最大的一处"地毯"，地毯下的第一块砖是第 22 章的欧拉乘积，第二块是第 24 章的 $\zeta$ 复分析。$\mathrm{li}(x)$ 在 $x\to0^+$ 有奇点，标准做法从 $2$（或"欧拉主值"从 $0$）积起，本书按教材惯例取 $\int_2^x$；表格中 $\mathrm{li}$ 的数值是标准计算值。李特尔伍德"无限次变号"的证明与斯丘斯数的确切上界属专业结果，本书只作观光。
\end{chengshi}

仪表盘读数已定。但流量图背后那座反应堆的内部——$\zeta$ 函数、它的零点、那条神秘的临界线——还藏着数学界悬赏一百万美元的问题。观光通道，下一章打开。
